\section{Data and Model}
\label{sec:data}
We test two different astrophysical datasets, representing different domains and physical properties. The first dataset is a \emph{stellar} multimodal dataset comprising 190,091 stars, observed by $5$ different modalities coming from different telescopes - \emph{Kepler}\citep{kepler:mission}, TESS\citep{tess:mission}, LAMOST\citep{lamost:survey}, APOGEE\citep{apogee:mission}, and Gaia\citep{gaia:mission}. Each modality is processed through a specific encoder, and the modalities are aligned with a model that factorizes the latent space into shared and specific subspaces. The result is one shared feature and $5$ per-modality features, each projected to the token embedding space of the LLM as its own small group of tokens. The second dataset is galaxies from the COSMOS-Web survey \citep{Casey2023_cosmos-web}, including two modalities: images and star-formation history (SFH). Similar to the stellar dataset, we use a pretrained model that aligns the two modalities (here using CLIP \cite{Radford2021_CLIP}) and project the frozen features into the LLM's token embedding. Notably, the two datasets are different in many aspects: number and type of modalities, pretrained models, and scientific questions. Using the two datasets in the same pipeline highlights general vs data-specific conclusions.\\  
\textbf{training prose}. One of the challenges in this pipeline is creating the training dataset. This should include ground-truth prose per object. Contrary to other fields, in astrophysics data is usually stored only in numerical catalogs, so one needs to create the prose from scratch. Our main method is a \emph{template-based} method: a deterministic program that picks a subset of facts from the catalog, chooses a phrasing from a bank of templates, and
writes prose. Every sentence is checked against the record by a verifier before it enters the corpus. The
check is per-fact: it guarantees that the facts a description states are true of the object, not
that the description mentions every unusual object in a group. Both datasets share two types of prose: description of a single object (single), and comparison between pairs (pair). We also test deeper reasoning capabilities with two additional tasks. In the stellar dataset we use population prose: a group of $N$ objects with description of the typical members, and a callout
of the one that does not belong and on which property. In the galaxy dataset we are interested in non trivial cross modal correlation and check if the model highlight those in the prose. In~\Cref{sec:results} we compare the template-based method with a second method: an \emph{LLM-based} generation, where we let DeepSeekV4 write the population prose from catalog facts, and verify their correctness. \todo{here we should elaborate slightly more - how the raw data looks like, what properties, how we create the training prose (random selsection of features, different styles, etc.). we shold also have visual example of raw data and resulting treaining prose in at least two different styles. this example should be neat, clean and visually clear}\\
\textbf{Input tokens.} To improve the diversity of the data and reduce overfitting, we apply several augmentations to the input. First, we add a random number of sentences from general astrophysics papers (from arXiv) and insert them around the data tokens. Second, we add label smoothing and a cosine learning-rate schedule. These augmentations apply to the single and pair records; population records are presented as the data tokens alone. \todo{add ablation table/figure of the arxiv cosine and diversity in appendix. They should not be seperated and treated as one augmentation, as we discussed in the past}


\begin{figure}[h]
  \centering
  % =====================================================================
% Talking Latents — the pipeline.  \input by the paper.  HORIZONTAL.
%
% General on purpose: no survey, no modality, no base-model, no track name in
% the STRUCTURE. The example strings are necessarily from one track (stellar);
% the caption says so.
%
% Trained vs frozen is shown with icons, in the "fire / ice" convention:
%   flame     = parameters updated by the prose loss
%   snowflake = parameters never updated here
%   multimodal features  ICE   (frozen encoders + frozen alignment upstream)
%   adapter network      FIRE
%   LLM                  ICE, with a small LoRA tag marked FIRE
% (The new bracket / token embeddings are trained too; they are not a
% computation block, so they carry no icon -- the caption says it.)
%
% Both example texts are VERBATIM, not invented:
%   question -> configs/stellar/filler_queries.yaml : single_target[0]
%   prose    -> cache/stellar_gen_eval/generations/astrosage8b.jsonl, a single-
%               mode generation at T=0.6, quoted in full. NB greedy (T=0) output
%               from this arm is a bullet list, not free prose -- the free-form
%               register appears once the template attractor is broken. State
%               the temperature if the figure is ever queried on this.
% Prompt order matches training_prompt_builder.py:_build_filler_before:
%   <BOS> <question> [OBJECTS] <emb...> [/OBJECTS] <prose> <EOS>
%
% Scope note: single/pair records and the pair bias exist on BOTH tracks; they
% are omitted for simplicity. The head is omitted here too -- this is the base
% pipeline; the head enters in Sec. 3.
%
% Icons are drawn in TikZ (no icon font in this TinyTeX). Previous vertical
% version: paper/archive/architecture_body_vertical_2026-09-06.tex.
% NB body file, not standalone: TinyTeX here lacks grfext.sty.
% =====================================================================
\definecolor{ice}{HTML}{3F87C4}
\definecolor{flameyellow}{HTML}{FFD166}
\resizebox{\textwidth}{!}{%
\begin{tikzpicture}[
  font=\small,
  >={Stealth[length=2.6mm,width=2.0mm]},
  box/.style   ={draw=deepteal, fill=palegrey, line width=1.0pt,
                 rounded corners=2.5pt, align=center, inner sep=5pt,
                 minimum height=13mm},
  tok/.style   ={draw=deepteal, fill=paleteal, line width=0.8pt,
                 rounded corners=1pt, minimum size=5.0mm, inner sep=0pt},
  brak/.style  ={font=\scriptsize\ttfamily, text=deepteal},
  ex/.style    ={font=\scriptsize\itshape, text=deepteal},
  dnote/.style ={font=\scriptsize, text=neutral},
  arr/.style   ={->, line width=1.0pt, draw=deepteal},
  % ---- icons -------------------------------------------------------
  pics/snow/.style={code={
    \fill[white] (0,0) circle (0.215);
    \foreach \a in {0,60,...,300}{
      \draw[ice, line width=0.55pt, line cap=round] (0,0) -- (\a:0.17);
      \draw[ice, line width=0.45pt, line cap=round]
        (\a:0.105) -- ++(\a+50:0.06)  (\a:0.105) -- ++(\a-50:0.06);
    }}},
  pics/flame/.style={code={
    \fill[white] (0,0.02) circle (0.215);
    \fill[coral]
      (0,-0.17) .. controls (-0.21,-0.14) and (-0.19,0.10) .. (-0.06,0.19)
                .. controls (-0.03,0.11) and (0.02,0.08) .. (0.03,0.02)
                .. controls (0.11,0.07) and (0.15,0.14) .. (0.13,0.21)
                .. controls (0.24,0.10) and (0.22,-0.14) .. (0,-0.17) -- cycle;
    \fill[flameyellow]
      (0,-0.16) .. controls (-0.09,-0.12) and (-0.08,0.00) .. (-0.01,0.05)
                .. controls (0.03,0.00) and (0.07,0.00) .. (0.09,-0.03)
                .. controls (0.11,-0.10) and (0.07,-0.15) .. (0,-0.16) -- cycle;
  }},
]

% ------------------------------------------------ input block (the anchor)
% token row chained with relative spacing so it survives any document font size
\node[brak] (ob) at (0,0) {[OBJECTS]};
\node[tok, right=1.8mm of ob] (g0) {};
\node[tok, right=1.2mm of g0] (g1) {};
\node[text=deepteal, right=1.0mm of g1, inner sep=0pt] (dots) {$\cdots$};
\node[tok, right=1.0mm of dots] (k0) {};
\node[tok, right=1.2mm of k0] (k1) {};
\node[brak, right=1.8mm of k1] (cb) {[/OBJECTS]};
\coordinate (rowmid) at ($(ob.west)!0.5!(cb.east)$);
\node[ex] (q) at ($(rowmid)+(0,0.50)$) {``Describe this star.''};
\begin{scope}[on background layer]
  \node[draw=deepteal, line width=1.0pt, rounded corners=2.5pt, fill=white,
        fit=(q)(ob)(cb), inner sep=4pt] (input) {};
\end{scope}
\node[dnote, anchor=north] at ($(input.south)+(0,-1mm)$) {\textbf{input}};

% ------------------------------------------------ upstream (left)
\node[box, left=7mm of input] (adapt) {\textbf{adapter}\\\textbf{network}};
\node[box, left=7mm of adapt]  (feat)
  {\textbf{multimodal}\\\textbf{features}\\[-1pt] \scriptsize $N$ objects};
\draw[arr] (feat) -- (adapt);
\draw[arr] (adapt) -- (input);

% ------------------------------------------------ LLM (right)
\node[box, right=7mm of input, minimum width=21mm] (llm) {\textbf{LLM}\\[3pt]
  {\scriptsize\ttfamily\color{deepteal} LoRA\hspace{5mm}}};
\draw[arr] (input) -- (llm);

% ------------------------------------------------ output prose (right)
\node[draw=deepteal, line width=1.0pt, rounded corners=2.5pt, fill=white,
      right=7mm of llm, align=left, inner sep=4pt, font=\scriptsize,
      text width=54mm]
  (prose)
  {This source presents as a catalogued main-sequence dwarf. Catalog values:
   spectroscopic Teff 6485 K; stellar log g 4.07 $(+0.03/{-}0.04)$; reported age
   (derived from isochrone) 2.1 Gyr. This star's [Fe/H] falls in the
   above-median quartile of the Kepler field.};
\node[dnote, anchor=north] at ($(prose.south)+(0,-1mm)$) {\textbf{prose}};
\draw[arr] (llm) -- (prose);

% ------------------------------------------------ icons: trained / frozen
\pic at ($(feat.north east)+(-0.6mm,-0.6mm)$)  {snow};
\pic at ($(adapt.north east)+(-0.6mm,-1.0mm)$) {flame};
\pic at ($(llm.north east)+(-0.6mm,-0.6mm)$)   {snow};
\pic at ($(llm.south)+(4.6mm,3.6mm)$)          {flame};

% ------------------------------------------------ legend
\node[dnote, anchor=north west] (leg) at ($(feat.south west)+(0,-2.0mm)$) {};
\pic at ($(leg.north west)+(2.0mm,-1.8mm)$) {flame};
\node[dnote, anchor=west] at ($(leg.north west)+(5.0mm,-1.8mm)$) {trained};
\pic at ($(leg.north west)+(21.0mm,-1.8mm)$) {snow};
\node[dnote, anchor=west] at ($(leg.north west)+(24.0mm,-1.8mm)$) {frozen};

\end{tikzpicture}}

 \caption{\textbf{The pipeline.} An adapter projects each object's multimodal
  features into tokens placed between brackets; the
  model continues from them into prose, and the loss is taken on the prose
  alone. The language model is adapted either by
  full fine-tuning or with LoRA.}
  \label{fig:arch}
\end{figure}
